% GENERATED FILE. Edit canonical lessons, then run npm run generate:books.
% canonical-source-hash: fnv1a64:f5d8ddc8239421cf
% canonical-lessons: MW-C203-atro, MW-C203-dag, MW-C203-disa, MW-C203-bhatakno, MW-C203-disno, MW-R203-review-198-203

\chapter{How Far Is It?}
\label{ch:mw-how-far-is-it}
\hlchaptermodality{203}

\begin{chapteropening}
\textbf{By the end of this chapter:} \emph{I can say how far away a place is.}

Complete the last lesson of chapter 203: I can say how far away a place is.
\end{chapteropening}

\section[\texorpdfstring{\mw{आंतरो} (ā̃tro)}{आंतरो (ā̃tro)}]{\mw{आंतरो} (ā̃tro) — a distance}

\label{lesson:MW-C203-atro}



\begin{quote}
Before the new one: say the Marwadi for crosswise, then the Marwadi for slanting.
\end{quote}

\subsection*{You'll want to know: \mw{आंतरो}}
\textbf{\mw{आंतरो}} — \emph{ā̃tro} — \textquotedblleft{}a distance\textquotedblright{}.

\begin{quote}
\mw{अठै} \mw{सूं} \mw{टेसण} \mw{तांई} \mw{घणो} \mw{आंतरो} \mw{है।}
\end{quote}

\begin{quote}
\emph{aṭhai sū̃ ṭesaṇ tā̃ī ghaṇo ā̃tro hai.}
\end{quote}

\begin{quote}
It's a long way from here to the station.
\end{quote}

A masculine noun.

\begin{cousinweb}
\textbf{\mw{आंतरो}} (\emph{ā̃tro}) is inherited from Sanskrit \emph{antara}, \textquotedblleft{}interior; an interval\textquotedblright{}. CDIAL gives no Marwari form in this entry, but lists Gujarati \emph{ā̃tro}, \textquotedblleft{}distance, interval\textquotedblright{} (CDIAL 357).
\end{cousinweb}

\begin{grammarlens}[title={What you've built}]
One more word for this chapter.
\end{grammarlens}

\subsection*{Your turn}
\begin{itemize}
\raggedright
  \item \emph{Say it:} \emph{ā̃tro}
  \item \emph{Say it:} \emph{ā̃tro}, once more
  \item \emph{Say it:} \emph{aṭhai sū̃ ṭesaṇ tā̃ī ghaṇo ā̃tro hai.} — It's a long way from here to the station.
  \item \emph{Recall:} say the Marwadi for crosswise, then the Marwadi for slanting, then say \emph{ā̃tro} again
\end{itemize}

\begin{tcolorbox}[breakable,colback=teal!4,colframe=teal!35!black,title={Before you move on}]
What does \mw{आंतरो} mean? (\textquotedblleft{}A distance\textquotedblright{}.) Say it once more.
\end{tcolorbox}
\section[\texorpdfstring{\mw{डग} (ḍag)}{डग (ḍag)}]{\mw{डग} (ḍag) — a step; a pace}

\label{lesson:MW-C203-dag}



\begin{quote}
Before the new one: say the Marwadi for slanting, then the Marwadi for a distance.
\end{quote}

\subsection*{You'll want to know: \mw{डग}}
\textbf{\mw{डग}} — \emph{ḍag} — \textquotedblleft{}a step; a pace\textquotedblright{}.

\begin{quote}
\mw{मंदिर} \mw{दस} \mw{डग} \mw{आगै} \mw{है।}
\end{quote}

\begin{quote}
\emph{mandir das ḍag āgai hai.}
\end{quote}

\begin{quote}
The temple is ten paces ahead.
\end{quote}

A feminine noun.

\begin{cousinweb}
\textbf{\mw{डग}} (\emph{ḍag}) is inherited from a reconstructed \emph{*ḍag}, \textquotedblleft{}step, pace\textquotedblright{}; CDIAL lists Old Marwari \emph{ḍaga}, \textquotedblleft{}step, stride\textquotedblright{} (CDIAL 5523).
\end{cousinweb}

\begin{grammarlens}[title={What you've built}]
One more word for this chapter.
\end{grammarlens}

\subsection*{Your turn}
\begin{itemize}
\raggedright
  \item \emph{Say it:} \emph{ḍag}
  \item \emph{Say it:} \emph{ḍag}, once more
  \item \emph{Say it:} \emph{mandir das ḍag āgai hai.} — The temple is ten paces ahead.
  \item \emph{Recall:} say the Marwadi for slanting, then the Marwadi for a distance, then say \emph{ḍag} again
\end{itemize}

\begin{tcolorbox}[breakable,colback=teal!4,colframe=teal!35!black,title={Before you move on}]
What does \mw{डग} mean? (\textquotedblleft{}A step; a pace\textquotedblright{}.) Say it once more.
\end{tcolorbox}
\section[\texorpdfstring{\mw{दिसा} (disā)}{दिसा (disā)}]{\mw{दिसा} (disā) — a direction}

\label{lesson:MW-C203-disa}



\begin{quote}
Before the new one: say the Marwadi for a distance, then the Marwadi for a pace.
\end{quote}

\subsection*{You'll want to know: \mw{दिसा}}
\textbf{\mw{दिसा}} — \emph{disā} — \textquotedblleft{}a direction\textquotedblright{}.

\begin{quote}
\mw{टेसण} \mw{किसी} \mw{दिसा} \mw{मांय} \mw{है}?
\end{quote}

\begin{quote}
\emph{ṭesaṇ kisī disā mā̃y hai?}
\end{quote}

\begin{quote}
Which way is the station?
\end{quote}

A feminine noun.

\begin{cousinweb}
\textbf{\mw{दिसा}} (\emph{disā}) is inherited from Sanskrit \emph{diśā}, \textquotedblleft{}direction, region\textquotedblright{}. CDIAL gives no Marwari form in this entry, but lists Hindi \emph{dis}, \textquotedblleft{}direction, side\textquotedblright{} (CDIAL 6339).
\end{cousinweb}

\begin{grammarlens}[title={What you've built}]
One more word for this chapter.
\end{grammarlens}

\subsection*{Your turn}
\begin{itemize}
\raggedright
  \item \emph{Say it:} \emph{disā}
  \item \emph{Say it:} \emph{disā}, once more
  \item \emph{Say it:} \emph{ṭesaṇ kisī disā mā̃y hai?} — Which way is the station?
  \item \emph{Recall:} say the Marwadi for a distance, then the Marwadi for a pace, then say \emph{disā} again
\end{itemize}

\begin{tcolorbox}[breakable,colback=teal!4,colframe=teal!35!black,title={Before you move on}]
What does \mw{दिसा} mean? (\textquotedblleft{}A direction\textquotedblright{}.) Say it once more.
\end{tcolorbox}
\section[\texorpdfstring{\mw{भटकणो} (bhaṭakṇo)}{भटकणो (bhaṭakṇo)}]{\mw{भटकणो} (bhaṭakṇo) — to lose one's way; to wander}

\label{lesson:MW-C203-bhatakno}



\begin{quote}
Before the new one: say the Marwadi for a pace, then the Marwadi for a direction.
\end{quote}

\subsection*{You'll want to know: \mw{भटकणो}}
\textbf{\mw{भटकणो}} — \emph{bhaṭakṇo} — \textquotedblleft{}to lose one's way; to wander\textquotedblright{}.

\begin{quote}
\mw{हूं} \mw{सहर} \mw{मांय} \mw{भटकूं} \mw{हूं।}
\end{quote}

\begin{quote}
\emph{hū̃ sahar mā̃y bhaṭkū̃ hū̃.}
\end{quote}

\begin{quote}
I keep getting lost in the city.
\end{quote}

Past (I …): \mw{हूं} \mw{भटक्यो}, \emph{hū̃ bhaṭakyo}.

\begin{cousinweb}
\textbf{\mw{भटकणो}} (\emph{bhaṭakṇo}) is inherited from Sanskrit \emph{bhraṣṭa}, \textquotedblleft{}fallen\textquotedblright{}. CDIAL gives no Marwari form in this entry, but lists Hindi \emph{bhaṭaknā}, \textquotedblleft{}to wander\textquotedblright{} (CDIAL 9655).
\end{cousinweb}

\begin{grammarlens}[title={What you've built}]
One more word for this chapter.
\end{grammarlens}

\subsection*{Your turn}
\begin{itemize}
\raggedright
  \item \emph{Say it:} \emph{bhaṭakṇo}
  \item \emph{Say it:} \emph{bhaṭakṇo}, once more
  \item \emph{Say it:} \emph{hū̃ sahar mā̃y bhaṭkū̃ hū̃.} — I keep getting lost in the city.
  \item \emph{Recall:} say the Marwadi for a pace, then the Marwadi for a direction, then say \emph{bhaṭakṇo} again
\end{itemize}

\begin{tcolorbox}[breakable,colback=teal!4,colframe=teal!35!black,title={Before you move on}]
What does \mw{भटकणो} mean? (\textquotedblleft{}To lose one's way; to wander\textquotedblright{}.) Say it once more.
\end{tcolorbox}
\section[\texorpdfstring{\mw{दीसणो} (dīsṇo)}{दीसणो (dīsṇo)}]{\mw{दीसणो} (dīsṇo) — to be visible; to appear}

\label{lesson:MW-C203-disno}



\begin{quote}
Before the new one: say the Marwadi for a direction, then the Marwadi for to lose one's way.
\end{quote}

\subsection*{You'll want to know: \mw{दीसणो}}
\textbf{\mw{दीसणो}} — \emph{dīsṇo} — \textquotedblleft{}to be visible; to appear\textquotedblright{}.

\begin{quote}
\mw{अठै} \mw{सूं} \mw{मंदिर} \mw{दीसै} \mw{है।}
\end{quote}

\begin{quote}
\emph{aṭhai sū̃ mandir dīsai hai.}
\end{quote}

\begin{quote}
The temple can be seen from here.
\end{quote}

Past (it …): \mw{वो} \mw{मंदिर} \mw{दीस्यो}, \emph{vo mandir dīsyo}.

\begin{cousinweb}
\textbf{\mw{दीसणो}} (\emph{dīsṇo}) is inherited from Sanskrit \emph{dṛśyate}, \textquotedblleft{}is seen\textquotedblright{}; CDIAL lists the Old Marwari verb \emph{dīsaï}, \textquotedblleft{}is seen\textquotedblright{} (CDIAL 6516).
\end{cousinweb}

\begin{grammarlens}[title={What you've built}]
One more word for this chapter.
\end{grammarlens}

\subsection*{Your turn}
\begin{itemize}
\raggedright
  \item \emph{Say it:} \emph{dīsṇo}
  \item \emph{Say it:} \emph{dīsṇo}, once more
  \item \emph{Say it:} \emph{aṭhai sū̃ mandir dīsai hai.} — The temple can be seen from here.
  \item \emph{Recall:} say the Marwadi for a direction, then the Marwadi for to lose one's way, then say \emph{dīsṇo} again
\end{itemize}

\begin{tcolorbox}[breakable,colback=teal!4,colframe=teal!35!black,title={Before you move on}]
What does \mw{दीसणो} mean? (\textquotedblleft{}To be visible; to appear\textquotedblright{}.) Say it once more.
\end{tcolorbox}
\section[(dialogue)]{Looking back: from At the Office to How Far Is It?}

\label{lesson:MW-R203-review-198-203}



\begin{quote}
Say the Marwadi for an office and an officer.
\end{quote}

\subsection*{Your turn}
Say these aloud:

\begin{itemize}
\raggedright
  \item an office, an officer, an application, a request, to explain
  \item an official stamp, a signature, a receipt, a reckoning, to write down
  \item scales, a seer, a kilo, a quarter, a litre
  \item a fort, a neighbourhood, a police station, an inn, to get in
  \item to turn, back again, the wrong way round, crosswise, slanting
  \item a distance, a pace, a direction, to lose one's way, to be visible
\end{itemize}

\begin{tcolorbox}[breakable,colback=teal!4,colframe=teal!35!black,title={Before you move on}]
An office? (\textbf{\mw{दफ्तर}}, \emph{daphtar}.) An officer? (\textbf{\mw{अफसर}}, \emph{aphsar}.) An application? (\textbf{\mw{अरजी}}, \emph{arjī}.) A request? (\textbf{\mw{बिनती}}, \emph{bintī}.) To explain? (\textbf{\mw{समझावणो}}, \emph{samjhāvṇo}.) An official stamp? (\textbf{\mw{मोहर}}, \emph{mohar}.) A signature? (\textbf{\mw{दसखत}}, \emph{daskhat}.) A receipt? (\textbf{\mw{रसीद}}, \emph{rasīd}.) A reckoning? (\textbf{\mw{लेखो}}, \emph{lekho}.) To write down? (\textbf{\mw{मांडणो}}, \emph{mā̃ḍṇo}.) Scales? (\textbf{\mw{तराजू}}, \emph{tarājū}.) A seer? (\textbf{\mw{सेर}}, \emph{ser}.) A kilo? (\textbf{\mw{किलो}}, \emph{kilo}.) A quarter? (\textbf{\mw{पाव}}, \emph{pāv}.) A litre? (\textbf{\mw{लीटर}}, \emph{līṭar}.) A fort? (\textbf{\mw{कोट}}, \emph{koṭ}.) A neighbourhood? (\textbf{\mw{मोहल्लो}}, \emph{mohallo}.) A police station? (\textbf{\mw{थाणो}}, \emph{thāṇo}.) An inn? (\textbf{\mw{सराय}}, \emph{sarāy}.) To get in? (\textbf{\mw{घुसणो}}, \emph{ghusṇo}.) To turn? (\textbf{\mw{मुड़णो}}, \emph{muṛṇo}.) Back again? (\textbf{\mw{पाछो}}, \emph{pāchho}.) The wrong way round? (\textbf{\mw{उलटो}}, \emph{ulṭo}.) Crosswise? (\textbf{\mw{आडो}}, \emph{āḍo}.) Slanting? (\textbf{\mw{तिरछो}}, \emph{tirchho}.) A distance? (\textbf{\mw{आंतरो}}, \emph{ā̃tro}.) A pace? (\textbf{\mw{डग}}, \emph{ḍag}.) A direction? (\textbf{\mw{दिसा}}, \emph{disā}.) To lose one's way? (\textbf{\mw{भटकणो}}, \emph{bhaṭakṇo}.) To be visible? (\textbf{\mw{दीसणो}}, \emph{dīsṇo}.)
\end{tcolorbox}
